% GPSS Theorem 2.3, complete constructive separation and its exact comparison.

\subsection{Separating past and future}
\label{sec:p09-gpss-past-future-separation}

GPSS Theorem~2.3 states that every \(L(S,U)\) formula is
equivalent, on discrete complete orders, to a Boolean
combination of \(L(S)\) and \(L(U)\) formulas. The
construction here passes through the complete first-order
normal form and the already proved future translation.
The paper separately refers to Gabbay's September 1979
manuscript for a direct syntactic proof. This supplied
construction proves the stated theorem without identifying
that unacquired manuscript with the present proof
\cite[Theorem~2.3 and its referral, p.~170; reference G2, p.~173]{GabbayPnueliShelahStavi1980}.

\begin{hbdefinition}{def:p09-gpss-separation-languages}{The two pure operator fragments}
Use the source's strict infix convention:
\(A\mathbin{\mathsf U_s}B\) requires \(B\) at a
strict future witness and \(A\) on the open gap;
\(A\mathbin{\mathsf S_s}B\) uses the corresponding
strict past witness. The language \(L(S,U)\)
contains proposition letters, Boolean constants,
Boolean operations, and those two constructors.
\(L(U)\) permits only the until constructor,
and \(L(S)\) only the since constructor.

Here \emph{pure} means that the other temporal
operator is absent. It does not exclude current-point
atoms: a Boolean combination of proposition letters
belongs to both fragments. A future formula may
depend on the present and future; a past formula
may depend on the present and past.

Derived abbreviations may first be expanded:
\[
\begin{aligned}
 FA&=\top\mathbin{\mathsf U_s}A,&
 PA&=\top\mathbin{\mathsf S_s}A,\\
 GA&=\neg F\neg A,&
 HA&=\neg P\neg A,\\
 XA&=\bot\mathbin{\mathsf U_s}A,&
 YA&=\bot\mathbin{\mathsf S_s}A.
\end{aligned}
\]
The last two use an immediate successor or predecessor,
respectively, if it exists. Empty open gaps retain
their ordinary vacuous universal meaning.
\end{hbdefinition}

\begin{hbdefinition}{def:p09-gpss-separation-translation}{The full first-order translation at a named point}
For a current variable \(x\), define \(A^*(x)\)
by translating \(q_i\) to \(q_i(x)\), commuting
with Boolean operations, and using fresh bound
variables in
\[
\begin{aligned}
 (A\mathbin{\mathsf U_s}B)^*(x)
 &=\exists y\left(x<y\land B^*(y)\land
                  \forall z(x<z<y\to A^*(z))\right),\\
 (A\mathbin{\mathsf S_s}B)^*(x)
 &=\exists y\left(y<x\land B^*(y)\land
                  \forall z(y<z<x\to A^*(z))\right).
\end{aligned}
\]
Each copied constituent receives fresh bound names;
only its free current-point variable is replaced.
Logical constants may be translated as \(x=x\)
and \(\neg(x=x)\).
\end{hbdefinition}

\begin{hblemma}{thm:p09-gpss-separation-translation}{proved=annex:p9.gpss.separation.translation}{The exact temporal-to-first-order bridge}
For every linear-order model, every point \(t\),
and every \(A\in L(S,U)\),
\[
                  \mathcal M,t\models A
                    \quad\Longleftrightarrow\quad
                  \mathcal M\models A^*(x)[x\mapsto t].
\]
The translation is finite and introduces no unary
predicate absent from the input. Its source argument
order agrees with Kamp's prefix notation by
\[
       U(B,A)=A\mathbin{\mathsf U_s}B,\qquad
       S(B,A)=A\mathbin{\mathsf S_s}B.
\]
Pure \(L(U)\) formulas depend only on the labelled
tail at and after the current point; pure \(L(S)\)
formulas depend only on the labelled initial segment
at and before it.
\end{hblemma}

\begin{hbdefinition}{def:p09-gpss-separation-dual}{Typed order reversal}
For a model \(\mathcal M\), let
\(\mathcal M^{\mathrm{op}}\) have the same domain
and unary interpretations, with the order reversed.
For a first-order formula \(\rho\), define
\(\rho^{\mathrm d}\) by reversing every order
atom, retaining equality and unary atoms, and
commuting with connectives and quantifiers.
For a temporal formula define
\[
\begin{aligned}
 (A\mathbin{\mathsf U_s}B)^{\mathrm d}
       &=A^{\mathrm d}\mathbin{\mathsf S_s}B^{\mathrm d},\\
 (A\mathbin{\mathsf S_s}B)^{\mathrm d}
       &=A^{\mathrm d}\mathbin{\mathsf U_s}B^{\mathrm d},
\end{aligned}
\]
with atoms and constants fixed and Boolean operations
retained. Reversal changes the time direction,
not the interval/witness operand positions.
Both operations are involutions, up to fresh
bound-variable names and the standard spelling
of order guards.
\end{hbdefinition}

\begin{hblemma}{thm:p09-gpss-separation-dual}{proved=annex:p9.gpss.separation.dual}{Reversal preserves the exact model and regional classes}
For a first-order formula \(\rho\) and assignment \(\nu\),
\[
           \mathcal M,\nu\models\rho^{\mathrm d}
                \quad\Longleftrightarrow\quad
           \mathcal M^{\mathrm{op}},\nu\models\rho.
\]
For a temporal formula \(A\) and actual point \(t\),
\[
           \mathcal M,t\models A^{\mathrm d}
                \quad\Longleftrightarrow\quad
           \mathcal M^{\mathrm{op}},t\models A.
\]
The first-order translation commutes with
these two reversals up to bound renaming.
Linearity, discreteness and Dedekind completeness
are preserved; first and last points, and
successors and predecessors, exchange roles.

On the exact GPSS regional stages with actual
endpoint variables \(a,b\), reversing
\(\operatorname{Bet}_{G,n}(a,b)\) yields
\(\operatorname{Bet}_{G,n}(b,a)\), after
standardizing the reversed guard. Consequently
\(\operatorname{Bef}_{G,n}(x)\) and
\(\operatorname{Aft}_{G,n}(x)\) exchange roles.
Formal ray endpoints exchange only as notation;
they never become domain points or predicate
arguments. Every before component is past,
and every after component is future.
\end{hblemma}

\begin{hbdefinition}{def:p09-gpss-separation-algorithm}{The effective discrete-complete separation algorithm}
Let \(T_F^{\mathrm{dc}}\) be the actual
finite future translation supplied by
\cref{thm:p09-gpss-discrete-future-completeness}.
For a before component define
\[
        T_P^{\mathrm{dc}}(\beta)
                   =\bigl(T_F^{\mathrm{dc}}
                                      (\beta^{\mathrm d})\bigr)^{\mathrm d}.
\]
The inner reversal is first-order; the outer
one is temporal. Its input is an after
component and its output belongs to \(L(S)\).

For \(A\in L(S,U)\), perform these finite steps:
\begin{enumerate}
\item Form \(A^*(x)\).
\item Apply \cref{thm:p09-gpss-special-normal-form}
      with declared parameter list \((x)\).
      Group each resulting clause as
      \[
       \bigwedge_j\beta_{ij}(x)\land
                     C_i(x)\land
       \bigwedge_k\gamma_{ik}(x),
      \]
      where the \(\beta_{ij}\) are before
      components, \(C_i\) is a conjunction
      of current signed literals, and the
      \(\gamma_{ik}\) are after components.
      The unique one-point order prefix \(x=x\)
      is a tautology.
\item Return
      \[
       \operatorname{Sep}_{\mathrm{dc}}(A)=
       \bigvee_i\left[
        \bigwedge_j T_P^{\mathrm{dc}}(\beta_{ij})
        \land C_i
        \land\bigwedge_k T_F^{\mathrm{dc}}(\gamma_{ik})
                    \right].
      \]
\end{enumerate}
Each empty conjunction is true and the empty
outer disjunction is false. The current literal
formula uses the same proposition letters as
its first-order clause.
\end{hbdefinition}

\begin{hbtheorem}{thm:p09-gpss-past-future-separation}{proved=annex:p9.gpss.past.future.separation}{GPSS Theorem 2.3 with its complete construction}
For every \(A\in L(S,U)\), the displayed
algorithm produces a Boolean combination
of pure \(L(S)\) and pure \(L(U)\) formulas
equivalent to \(A\) at every point of every
nonempty discrete complete model, for every
unary interpretation.

This includes finite, natural and integer
orders, endpoint points and empty past or
future regions. It is not restricted to
the initial point of an \(\omega\)-model.
All calls are to supplied finite constructions;
no search for a semantically equivalent formula
is part of the algorithm
\cite[Theorem~2.3, p.~170]{GabbayPnueliShelahStavi1980}.
\end{hbtheorem}

\begin{hbdefinition}{def:p09-gpss-guarded-separation-bases}{The two guarded bases for complete-order separation}
Let \(\mathcal A_F=TL(U,K^-)\) with the
vacuous past-limit convention already fixed.
Define its dual limit by
\[
 K^+_{\mathrm v}B(t)=
       \forall r\left(t<r\to
                    \exists u(t<u<r\land B(u))\right)
       =\neg\bigl((\neg B)\mathbin{\mathsf U_s}\top\bigr)(t).
\]
The subscript distinguishes this vacuous
right limit from a nonempty right limit.
It is true at a last point.
Put \(\mathcal A_P=TL(S,K^+_{\mathrm v})\).

Let \(\mathcal F\) be the Boolean closure
of atoms, constants and \(B\mathbin{\mathsf U_s}C\)
with \(B,C\in\mathcal A_F\); let \(\mathcal P\)
be the corresponding closure using
\(B\mathbin{\mathsf S_s}C\) with
\(B,C\in\mathcal A_P\).
Equivalently every \(K^-\) in \(\mathcal F\)
lies under until, and every \(K^+_{\mathrm v}\)
in \(\mathcal P\) lies under since.

Extend temporal reversal by
\[
               (K^-B)^{\mathrm d}
                    =K^+_{\mathrm v}(B^{\mathrm d}),\qquad
               (K^+_{\mathrm v}B)^{\mathrm d}
                    =K^-(B^{\mathrm d}).
\]
The guarded future translation from
\cref{thm:p09-guarded-future-completeness}
is denoted \(T_F^{\mathrm{com}}\).
For a before component use
\[
        T_P^{\mathrm{com}}(\beta)=
           \bigl(T_F^{\mathrm{com}}
                            (\beta^{\mathrm d})\bigr)^{\mathrm d}.
\]
\end{hbdefinition}

\begin{hblemma}{thm:p09-gpss-guarded-separation-dual}{proved=annex:p9.gpss.guarded.separation.dual}{The guarded past translation and its semantic scope}
The extended reversal retains its satisfaction
identity and maps \(\mathcal F\) onto
\(\mathcal P\). Every formula in
\(\mathcal F\) is future and every formula
in \(\mathcal P\) is past, on all linear
orders. On complete orders,
\(T_P^{\mathrm{com}}(\beta)\) expresses its
before component at every point.

Each guarded formula expands finitely into
\(L(S,U)\), using the displayed right-limit
clause and
\(K^-B=\neg((\neg B)\mathbin{\mathsf S_s}\top)\).
Thus the guarded comparison introduces no
undefined temporal operation.
\end{hblemma}

\begin{hbtheorem}{thm:p09-gpss-complete-order-separation}{proved=annex:p9.gpss.complete.order.separation}{The correct complete-order comparison}
Replacing the two discrete translators in
\(\operatorname{Sep}_{\mathrm{dc}}\) by
\(T_F^{\mathrm{com}},T_P^{\mathrm{com}}\)
gives an effective separation of every
\(L(S,U)\) formula into a Boolean combination
of \(\mathcal F\) and \(\mathcal P\) formulas
over all nonempty Dedekind-complete orders.
The components are semantically future and
past, with current atoms permitted.

Their internal limit operators need not
belong to the pure \(L(U)\) or \(L(S)\)
grammar. The result uses the exact guarded
completeness theorem and its order dual;
it does not extend GPSS's pure-operator
conclusion unchanged to all complete orders
\cite[Theorem~7.4]{Rabinovich2014Kamp}.
\end{hbtheorem}

\begin{hbtheorem}{thm:p09-gpss-pure-separation-obstruction}{proved=annex:p9.gpss.pure.separation.obstruction}{Pure since/until separation fails over the real line}
The mixed formula
\[
 \Xi=
 \top\mathbin{\mathsf U_s}
       \left(P\land
          \neg\bigl((\neg P)\mathbin{\mathsf S_s}\top\bigr)\right)
\]
has no equivalent Boolean combination of
pure \(L(S)\) and pure \(L(U)\) formulas
over all complete orders, and already
has no such equivalent over the real line.

Use the models \(C_1,C_2\) from
\cref{def:p09-hr-continuous-models}, with
their block beginnings marked. At their
least points every pure-since formula
agrees, and every pure-until formula
agrees by the complete four-truth-set
calculation of \cref{def:p09-hr-until-table}.
The formula \(\Xi\) is false in \(C_1\)
and true in \(C_2\). Extending both
models by the same unmarked negative
real half-line proves the full-real-line
claim while preserving the pure-past
and pure-future agreements.
\end{hbtheorem}

\begin{hbexample}{def:p09-gpss-separation-worked}{A future point at which something is already past}
On every linear order,
\[
 F(Pq)\longleftrightarrow
       \bigl((Pq\lor q)\land F\top\bigr)
                    \lor F(q\land F\top).
\]
Here \(P\) is the strict-past modality
in the derived abbreviation \(Pq\);
\(q\) is a proposition letter.
The right side is a Boolean combination
of a pure-past formula, current atoms
and pure-future formulas.

On the two-point order \(0<1\), with
\(q\) true only at 1, both sides are
false at zero. The future occurrence
of \(q\) needs a still later point
to count as past there. At a last
point the entire expression is false,
even if \(q\) has occurred in its past.
The annex proves the identity by the
three exact positions of the past witness.
\end{hbexample}
